\subsection{函数的限制与扩展}

两个函数f和g被称为在集合E上一致（或重合），如果E包含在f和g的定义域中，并且如果 \(f(x) = g(x)\) 对于所有 \(x \in E\) 两个具有相同图像的函数在其定义域上一致。说 f = g 就是说 f 和 g 具有相同的定义域 A 和相同的目标集 B，并且它们在 A 上一致。

设 \(f = (\mathbf{F}, \mathbf{A}, \mathbf{B})\) 和 \(g = (\mathbf{G}, \mathbf{C}, \mathbf{D})\) 是两个函数。说 \(\mathbf{F} \subset \mathbf{G}\) ，就是指 \(\mathbf{A}\) （ \(f\) 的定义域）包含在 \(\mathbf{C}\) （ \(g\) 的定义域）中，且 \(g\) 与 \(f\) 在 \(\mathbf{A}\) 上一致。如果还有 \(\mathbf{B} \subset \mathbf{D}\) ，则称 \(g\) 是 \(f\) 的一个扩张（更准确地说，是 \(f\) 到 \(\mathbf{C}\) 上的扩张），也称 \(g\) 扩张了 \(f\) （到 \(\mathbf{C}\) 上）。当把 \(g\) 称为 \(\mathbf{D}\) 中的元素族时，就称 \(f\) 为 \(g\) 的子族。

设 \(f\) 是一个函数，且 \(A\) 是 \(f\) 的定义域的一个子集。显然，关系“ \(x \in A\) 且 \(y = f(x)\) ”关于 \(x\) 和 \(y\) 有一个图 \(G\) ，该图是函数图，且定义域为 \(A\) ；以 \(G\) 为图且与 \(f\) 有相同陪域的函数称为 f 在 A 上的限制，有时记作 \(f|A\) 。任一函数都是其每个限制的扩张。如果两个函数 f、g 有相同的陪域，并且在集合 E 上一致，则它们在 E 上的限制相等。

